\section{What is a quiver?}\label{sec:11-path-algebras:01-what-is-a-quiver:1}

\subsection{Recall}\label{sec:11-path-algebras:01-what-is-a-quiver:2}

Formal definition cited in this section, gathered here for quick reference (full citation in the \hyperref[sec:root:bibliography:1]{Bibliography}):

\begin{itemize}
\tightlist
\item
  \textbf{Quiver.} Stated verbatim as \(Q = (Q_0, Q_1, s, t)\) (\cite{AssemSimsonSkowronski2006}, Definition 1.1, Ch. II §1 ``Quivers and path algebras''): a set of vertices \(Q_0\), a set of arrows \(Q_1\), and two functions \(s, t : Q_1 \to Q_0\) giving each arrow's source and target.
\end{itemize}

A \textbf{quiver} is a \emph{directed graph}: a set of vertices and a set of directed edges (called \textbf{arrows}) between them. This is the same notion Mathlib calls \href{https://loogle.lean-lang.org/?q=Quiver}{\texttt{Quiver}} (built here from scratch, following Chapter 0's ``no Mathlib'' policy, rather than reusing Mathlib's). Formally, a quiver \(Q\) consists of:

\begin{itemize}
\tightlist
\item
  a set of vertices \(Q_0\),
\item
  a set of arrows \(Q_1\),
\item
  two functions \(s, t : Q_1 \to Q_0\) giving each arrow's \textbf{source} and \textbf{target}.
\end{itemize}

Picture an arrow \(\alpha : i \to j\) as a literal arrow drawn from vertex \(i\) to vertex \(j\); \(s(\alpha) = i\) and \(t(\alpha) = j\).

\subsection{Example quiver}\label{sec:11-path-algebras:01-what-is-a-quiver:3}

Take vertices \(\{1, 2, 3\}\) and arrows

\[
\alpha : 1 \to 2, \qquad \beta : 2 \to 3
\]

\subsection{References}\label{sec:11-path-algebras:01-what-is-a-quiver:4}

Full citations in the \hyperref[sec:root:bibliography:1]{Bibliography}. Formal definitions are gathered in Recall, above.

\begin{itemize}
\tightlist
\item
  Assem, Simson, and Skowroński (\cite{AssemSimsonSkowronski2006}), Definition 1.1, Ch. II §1 ``Quivers and path algebras'' --- quiver.
\item
  Schiffler (\cite{Schiffler2014}) --- an elementary, textbook-level treatment of the same definition.
\end{itemize}
